% !TeX program = xelatex
\documentclass[11pt,a4paper]{article}
\usepackage[utf8]{inputenc}
\usepackage[T1]{fontenc}
\usepackage[english]{babel}
\usepackage{amsmath,amssymb,amsfonts,mathtools,amsthm}
\usepackage{geometry}
\geometry{left=1.25cm,right=1.25cm,top=1.25cm,bottom=3.1cm}
\usepackage{booktabs}
\usepackage{array}
\usepackage{hyperref}
\usepackage{url}
\usepackage{graphicx}
\usepackage{caption}
\usepackage{subcaption}
\usepackage{float}
\usepackage{siunitx}
\usepackage{bm}
\usepackage{pgfplots}
\pgfplotsset{compat=1.18}
\usepackage{xcolor}
\usepackage{titlesec}
\usepackage{fancyhdr}
\usepackage{microtype}
\usepackage{fontspec}
\usepackage{tcolorbox}
\usepackage{enumitem}
\usepackage{xeCJK}  % 日本語サポート

% ラテン文字フォント（オリジナルと同じ）
\setmainfont{Calibri}
\setsansfont{Segoe UI}[
BoldFont = Segoe UI Bold,
Ligatures = TeX
]

% 日本語フォント（Windows標準）
\setCJKmainfont{MS Gothic}
\setCJKsansfont{MS Gothic}
\setCJKmonofont{MS Gothic}

\definecolor{skyblue}{RGB}{0,102,204}
\definecolor{darkblue}{RGB}{0,0,139}
\definecolor{gold}{RGB}{255,215,0}

% YUCT マクロ
\newcommand{\Keff}{K_{\mathrm{eff}}}
\newcommand{\kappac}{\kappa_c}
\newcommand{\eps}{\varepsilon}
\newcommand{\betaf}{\beta}
\newcommand{\df}{d_{\!f}}
\newcommand{\Sodd}{S_{\mathrm{odd}}}
\newcommand{\Seven}{S_{\mathrm{even}}}
\newcommand{\Mpl}{M_{\mathrm{Pl}}}
\newcommand{\Lpl}{l_{\mathrm{Pl}}}
\newcommand{\tpl}{t_{\mathrm{Pl}}}
\newcommand{\Lmin}{\ell_{\min}}
\newcommand{\LUV}{\Lambda_{\mathrm{UV}}}

% 定理環境（日本語）
\newtheorem{theorem}{定理}[section]
\newtheorem{proposition}[theorem]{命題}
\newtheorem{definition}[theorem]{定義}
\theoremstyle{remark}
\newtheorem{remark}[theorem]{注記}

% セクション書式
\titleformat{\section}{\LARGE\bfseries\color{darkblue}}{\thesection}{1em}{}
\titleformat{\subsection}{\Large\bfseries\color{darkblue}}{\thesubsection}{1em}{}
\titleformat{\subsubsection}{\normalsize\bfseries\color{darkblue}}{\thesubsubsection}{1em}{}

% ヘッダー・フッター（日本語版）
\fancypagestyle{firstpage}{%
	\fancyhf{}%
	\renewcommand{\headrulewidth}{-5pt}%
	\renewcommand{\footrulewidth}{-5pt}%
	\fancyfoot[C]{%
		\begin{minipage}[t]{\textwidth}
			\centering
			\textcolor{gold}{\rule{\textwidth}{1.6pt}}\\[5pt]
			{\small \textsuperscript{1}Yakushev Research, \url{https://yuct.org/}} \hfill
			{\small YPSDC プロトコル: \url{https://ypsdc.com/}}\\[-2pt]
			\textcolor{gold}{\rule{\textwidth}{1.6pt}}\\[5pt]
			\textbf{ヤクシェフ統一調整理論 2026} \hfill \thepage\\[10pt]
		\end{minipage}%
	}%
}

\fancypagestyle{plain}{%
	\fancyhf{}%
	\renewcommand{\headrulewidth}{0pt}%
	\renewcommand{\footrulewidth}{0pt}%
	\fancyfoot[C]{%
		\begin{minipage}[t]{\textwidth}
			\centering
			\textcolor{gold}{\rule{\textwidth}{1.6pt}}\\[4pt]
			\textbf{ヤクシェフ統一調整理論 2026} \hfill \thepage\\[4pt]
		\end{minipage}%
	}%
}

\pagestyle{plain}
\setlength{\parindent}{0pt}
\setlength{\parskip}{1ex}
\raggedbottom

\hypersetup{
	colorlinks=true,
	linkcolor=blue,
	citecolor=blue,
	urlcolor=blue,
}

% カスタムヘッダー
\newcommand{\makeheader}{%
	\noindent
	\begin{minipage}{0.17\textwidth}
		\raggedright
		{\fontspec{Segoe UI}[BoldFont=Segoe UI Bold]\bfseries\color{skyblue}\fontsize{46}{46}\selectfont YUCT}%
	\end{minipage}%
	\hspace{30pt}
	\begin{minipage}{0.32\textwidth}
		\raggedright
		{\sffamily\fontsize{11}{13}\selectfont YAKUSHEV\\ UNIFIED\\ COORDINATION THEORY}%
	\end{minipage}%
	\hfill
	\begin{minipage}{0.38\textwidth}
		\raggedleft
		{\sffamily\bfseries 付録 UVD}\\ 
		{\sffamily バージョン 1.0 / 2026年6月}\\ 
		{\sffamily\footnotesize \url{https://doi.org/10.5281/zenodo.18444598}}%
	\end{minipage}
	\par\vspace{5pt}
	\noindent\textcolor{gold}{\rule{\textwidth}{1.6pt}}
	\par\vspace{0pt}
}

\begin{document}
	\thispagestyle{firstpage}
	\makeheader
	
	\vspace{0.5cm}
	{\LARGE\bfseries 付録 UVD：YUCT における紫外発散の解決 \\
		\Large 調整形状因子と調整ネットワークの物理的スケールに基づく有限量子場理論 \par}
	\vspace{0.3cm}
	
	{\large Alexey V. Yakushev\textsuperscript{1}} \\
	\vspace{0.2cm}
	
	{\color{darkblue}\textbf{\LARGE 要旨}} \par
	{\large
		\noindent
		ヤクシェフ統一調整理論（YUCT）は，量子場理論における紫外発散が，真空の有限な情報容量から生じる調整形状因子 \(F(q^2) = \exp[-(q^2/\LUV^2)^{2/3}]\) によって除去されることを提案する。元々の定式化では，紫外スケール \(\LUV\) はプランク質量 \(\LUV = 3\pi\Mpl \approx 1.15\times 10^{20}\) GeV と同一視され，ヒッグス質量に対する巨大な有限放射補正をもたらした。我々は，スカラー質量補正に現れる定数 \(C_0\) の計算誤差を修正する：正しい値は \(C_0 = 3\sqrt{\pi}/(8\sqrt{2}) \approx 0.46999\) であり，以前報告された \(C_0 \approx 0.626\) を置き換える。この修正により，プランクスケールネットワークからのヒッグス質量への放射寄与は \(\delta m \approx 1.57\times 10^{18}\) GeV となり，裸の質量が正確に相殺されない限り深刻な微調整問題を再導入する。我々は二つの論理的代替案を検討する：(i) YUCT 質量公式 \(m = \Mpl(2/3)^L\xi_f\) がすでに物理質量を決定しており，放射補正は有限の相殺項によって吸収される – 技術的には整合しているが美的に満足できない解決策；(ii) プランクスケールの基本ラグランジアンは裸の質量項を含まず（スケール不変性），ヒッグス質量全体が放射補正によって生成される。この第二のシナリオでは，自己無撞着性から紫外スケールが \(\LUV \approx 8.96\) TeV であることが要求される – LHC でアクセス可能な値である。調整形状因子はその後，TeV スケールでの高エネルギー散乱断面積に観測可能なずれを予言する。我々は両方のシナリオとその数学的正当性，実験的結果を示し，YUCT を厳密な実証的テストにさらすことを目的とする。
	}
	
	\vspace{0.2cm}
	\noindent
	{\color{darkblue}\textbf{キーワード:}} YUCT，紫外発散，調整形状因子，階層性問題，ヒッグス質量，スケール不変性，非局所量子場理論，LHC 現象論。
	\vspace{0.1cm}
	\noindent
	{\color{darkblue}\textbf{引用:}} Yakushev AV. 付録 UVD：YUCT における紫外発散の解決. 2026. DOI: \url{https://doi.org/10.5281/zenodo.18444598}（本巻）.
	
	\vspace{0.5cm}
	
	% ============================================================
	\section{序論}
	\label{sec:intro}
	
	ヤクシェフ統一調整理論（YUCT）は，量子場理論（QFT）の紫外発散に対するエレガントな解決法を提供する。任意の運動量カットオフや無限大の再正規化を導入する代わりに，YUCT は物理的真空が量子情報を伝送するための有限な「帯域幅」を持つと仮定する。この帯域幅は，特性スケール \(\LUV\) 以上の仮想運動量を指数関数的に抑制する \textbf{調整形状因子} \(F(q^2)\) によって定量化される。この形状因子の数学的構造は，YUCT の普遍誤差則 \(\varepsilon \propto K_{\mathrm{eff}}^{-\beta}\)（\(\beta = 2/3\)）によって決定される。
	
	この付録の以前のバージョンでは，\(\LUV\) をプランク質量 \(\LUV = 3\pi\Mpl \approx 1.15\times 10^{20}\) GeV と同一視した。これは，YUCT の一次代数的ループによって与えられる最小長 \(\Lmin = \frac{2}{3}\Lpl\) から導かれるスケールである。この同一視により，すべてのループ積分は有限になる：対数発散は \(\ln(\LUV/m)\) に置き換えられ，べき乗発散は完全に消える。
	
	しかし，スカラー質量補正の慎重な再計算は，有限な残差の大きさを支配する定数 \(C_0\) の計算誤差を明らかにする。\(C_0\) の正しい値は以前報告されたものより小さく，放射補正は減少するが，それでも観測されたヒッグス質量よりも多くの桁だけ大きいままである。これにより，物理スケール \(\LUV\) の批判的な再検討が強いられる。
	
	この付録には二つの目的がある。第一に，修正された数学的導出を完全な詳細とともに提示し，すべての式が正確で再現可能であることを保証する。第二に，修正された式の物理的含意を探求する。これには，調整ネットワークの基本スケールがプランクスケールではなく，はるかに低いスケールである可能性があり，現在の衝突型加速器実験で到達可能であるという急進的な提案を含む。
	
	% ============================================================
	\section{調整形状因子と紫外スケール}
	\label{sec:formfactor}
	
	YUCT 調整形状因子は次式で与えられる：
	\begin{equation}
		F(q^2) = \exp\!\Bigl[-\Bigl(\frac{q^2}{\LUV^2}\Bigr)^{\!\beta}\,\Bigr], \qquad \beta = \frac{2}{3}.
		\label{eq:F}
	\end{equation}
	すべての標準模型の伝播関数は \(F(p^2/\LUV^2)\) で乗じられる。スケール \(\LUV\) は元々，コンプトン関係 \(\LUV = 2\pi\hbar c/\Lmin = 3\pi\Mpl\) を通じて最小長 \(\Lmin = \frac{2}{3}\Lpl\) から導かれる。
	
	この形状因子は，すべてのループ積分を紫外領域で絶対収束させる。証明は簡単である：ループ運動量における任意の多項式 \(P(k)\) に対して，積 \(P(k) \cdot \prod_i F(k_i^2/\LUV^2)\) は，指数関数が支配的であるため，\(|k| \to \infty\) で \(|k|\) の任意の逆べき乗よりも速く減衰する。
	
	% ============================================================
	\section{電子自己エネルギー：対数発散の解決}
	\label{sec:electron}
	
	QED における１ループ電子自己エネルギーは，従来の QFT では対数発散する。YUCT 形状因子を用いると，ユークリッド積分は
	\begin{equation}
		\Sigma_{\mathrm{YUCT}} \propto \int_0^{\infty} \frac{dk_E}{k_E}\; \exp\!\Bigl[-2\Bigl(\frac{k_E^2}{\LUV^2}\Bigr)^{\!2/3}\,\Bigr].
		\label{eq:elec_int}
	\end{equation}
	となる。置換 \(t = (k_E^2/\LUV^2)^{2/3}\) により \(dk_E/k_E = \frac{3}{4} dt/t\) が得られ，積分は
	\begin{equation}
		I_{\mathrm{log}} = \frac{3}{4} \int_{t_{\min}}^{\infty} \frac{dt}{t}\, e^{-2t} \approx \ln\frac{\LUV}{m_e},
		\label{eq:elec_result}
	\end{equation}
	と評価される。ここで \(t_{\min} \sim (m_e^2/\LUV^2)^{2/3}\) である。対数発散は有限な対数 \(\ln(\LUV/m_e) \approx 53.8\) で置き換えられる。この計算は初等的であり，修正を必要としない。
	
	% ============================================================
	\section{スカラー質量補正：定数 \(C_0\) の修正}
	\label{sec:scalar}
	
	\subsection{積分とその評価}
	
	\(\phi^4\) 理論における１ループスカラー質量補正は，ウィック回転と形状因子の挿入後，
	\begin{equation}
		\delta m^2_{\mathrm{YUCT}} = \frac{\lambda}{16\pi^2} \int_0^{\infty} \frac{k_E^3\, dk_E}{k_E^2 + m^2}\; F^2\!\Bigl(\frac{k_E^2}{\LUV^2}\Bigr), \qquad F^2(x) = e^{-2x^{2/3}}.
		\label{eq:scalar_YUCT}
	\end{equation}
	となる。
	
	無次元変数 \(t = k_E^2/\LUV^2\) を導入し，\(m \ll \LUV\) で展開すると，主要項は
	\begin{equation}
		\delta m^2_{\mathrm{YUCT}} = \frac{\lambda}{32\pi^2}\,\LUV^2 \int_0^{\infty} e^{-2t^{2/3}} dt + \mathcal{O}(m^2/\LUV^2).
		\label{eq:scalar_lead}
	\end{equation}
	である。
	
	以前（バージョン 3.1）では積分 \(C_0 \equiv \int_0^{\infty} e^{-2t^{2/3}} dt \approx 0.626\) と報告されていた。この値は誤りである。ここで正確に評価する。
	
	\subsection{\(C_0\) の正確な評価}
	
	\begin{theorem}[\(C_0\) の正確な値]
		\label{thm:C0}
		\[
		C_0 = \int_0^{\infty} e^{-2t^{2/3}} dt = \frac{3\sqrt{\pi}}{8\sqrt{2}} \approx 0.46999.
		\]
	\end{theorem}
	
	\begin{proof}
		\(u = 2t^{2/3}\) と置換する。すると \(t = (u/2)^{3/2}\) であり，
		\[
		dt = \frac{3}{2} (u/2)^{1/2} \cdot \frac{1}{2} du = \frac{3}{4\sqrt{2}} u^{1/2} du.
		\]
		したがって，
		\[
		C_0 = \int_0^{\infty} e^{-u} \cdot \frac{3}{4\sqrt{2}} u^{1/2} du = \frac{3}{4\sqrt{2}} \int_0^{\infty} u^{1/2} e^{-u} du = \frac{3}{4\sqrt{2}} \, \Gamma(3/2).
		\]
		\(\Gamma(3/2) = \frac{\sqrt{\pi}}{2}\) を用いて，
		\[
		C_0 = \frac{3}{4\sqrt{2}} \cdot \frac{\sqrt{\pi}}{2} = \frac{3\sqrt{\pi}}{8\sqrt{2}} \approx 0.46999.
		\]
	\end{proof}
	
	\subsection{修正されたスカラー質量補正}
	
	正しい \(C_0\) の値を用いると，スカラー質量補正の主要項は
	\begin{equation}
		\boxed{\delta m^2_{\mathrm{YUCT}} = \frac{\lambda}{32\pi^2}\,\LUV^2 \cdot \frac{3\sqrt{\pi}}{8\sqrt{2}} + \mathcal{O}(m^2/\LUV^2)}.
		\label{eq:scalar_corrected}
	\end{equation}
	となる。
	
	ヒッグスボソンについては \(\lambda = m_H^2/(2v^2) \approx 0.13\)（ここで \(v \approx 246\) GeV）である。プランクスケール同一視 \(\LUV = 3\pi\Mpl \approx 1.15\times 10^{20}\) GeV を用いると：
	\begin{align}
		\delta m^2_H &\approx \frac{0.13}{32\pi^2} \cdot (1.15\times 10^{20})^2 \cdot 0.46999 \notag\\
		&\approx 2.47 \times 10^{36}\ \mathrm{GeV}^2, \notag\\
		\delta m_H &\approx 1.57 \times 10^{18}\ \mathrm{GeV}.
		\label{eq:planck_correction}
	\end{align}
	
	この値は観測されたヒッグス質量（\(m_H^{\mathrm{exp}} \approx 125\) GeV）を \(1.26 \times 10^{16}\) 倍も超えている。ラグランジアン中の裸の質量がゼロならば，物理質量は完全にこの放射補正によって決定され，予測値は明らかに観測と両立しない。
	
	% ============================================================
	\section{物理的含意：二つのシナリオ}
	\label{sec:scenarios}
	
	修正された計算は，階層性問題に直接対峙することを余儀なくさせる。我々は論理的に整合した二つの代替案を検討する。
	
	\subsection{シナリオ A：有限相殺項を持つプランクスケールネットワーク}
	
	このシナリオでは，YUCT は同一視 \(\LUV = 3\pi\Mpl\) を維持する。任意の粒子の物理質量は YUCT 質量公式
	\begin{equation}
		m_{\mathrm{phys}} = \Mpl \left(\frac{2}{3}\right)^L \xi_f,
		\label{eq:mass_formula}
	\end{equation}
	で与えられる。ここで \(L\) は調整深度，\(\xi_f\) は共鳴因子である（付録 J 参照）。ラグランジアン中の裸の質量は自由パラメータであり，放射補正を打ち消し観測質量を再現するように選ばれる。
	
	数学的にはこれは完全に整合している：理論は有限であり，必要な相殺項も有限であり，無限大に遭遇することは決してない。YUCT 質量公式 Eq.~(\ref{eq:mass_formula}) は，相殺項を固定する再正規化条件を与える。
	
	しかし物理的には，このシナリオは裸の質量と放射補正が \(10^{16}\) 分の１の精度で打ち消し合うことを要求する – これはまさに階層性問題が回避しようとしていた種類の微調整である。YUCT はなぜこの打ち消しを達成するために裸の質量を選ぶべきかを説明しない；それは単に通常の無限大の引き算を巨大な有限の引き算で置き換えるだけである。
	
	\subsection{シナリオ B：スケール不変な起源と TeV スケールネットワーク}
	
	代わりとなる，美的により魅力的な解決法は，\(\LUV\) のプランクスケール同一視を放棄することである。もしプランクスケールの基本理論が裸の質量項を持たない（すなわちスケール不変である）ならば，すべての質量は調整形状因子によって放射的に生成される。この場合，物理的なヒッグス質量 \emph{は} 放射補正そのものである：
	\begin{equation}
		m_H^2 = \delta m^2_{\mathrm{YUCT}} = \frac{\lambda}{32\pi^2}\,\LUV^2 \cdot C_0.
		\label{eq:scale_inv}
	\end{equation}
	
	\(\LUV\) について解くと：
	\begin{equation}
		\LUV = \sqrt{\frac{32\pi^2\, m_H^2}{\lambda\, C_0}}.
		\label{eq:LUV_solve}
	\end{equation}
	
	修正された \(C_0 = 3\sqrt{\pi}/(8\sqrt{2}) \approx 0.46999\)，\(\lambda = 0.13\)，\(m_H = 125\) GeV を代入すると，
	\begin{equation}
		\boxed{\LUV \approx 8.96\ \mathrm{TeV}}.
		\label{eq:LUV_TeV}
	\end{equation}
	を得る。
	
	この値は理論の予測であり，自由パラメータではない。これは基本調整スケールを TeV 範囲に置き，実験的にアクセス可能にする。
	
	\subsection{シナリオ B の物理的解釈}
	
	もし \(\LUV \approx 9\) TeV ならば，調整ネットワークはプランクスケールではなく電弱スケールで「結晶化」する。形状因子 \(F(q^2) = \exp[-(q^2/\LUV^2)^{2/3}]\) は数 TeV 以上の仮想運動量を抑制し始める。これは直ちに観測可能な結果をもたらす：
	
	\begin{enumerate}[label=(\arabic*)]
		\item \textbf{LHC におけるドレル–ヤンおよびダイジェット断面積} は，標準模型の予測と比較して，高い不変質量での事象の不足を示すべきである。この抑制は \(m_{jj} \gtrsim 5\) TeV で顕著になる。
		\item \textbf{ヒッグス対生成} およびベクトルボソン散乱は，形状因子がヒッグス自己結合を変更するため，異常なエネルギー依存性を示すはずである。
		\item \textbf{電弱精密観測量のずれ} が予想されるかもしれないが，LEP 測定は \(\LUV\) よりはるかに低いスケールを探査しているため，それらはおそらく小さい。
	\end{enumerate}
	
	決定的に，このシナリオでの最小長は \(\Lmin = 2\pi\hbar c/\LUV \approx 1.38\times 10^{-19}\) m であり，プランク長（\(\sim 10^{-35}\) m）よりもはるかに大きい。これは，時空の粗さが量子重力効果ではなく，電弱スケールの現象であり，おそらくヒッグス機構自体と関連していることを意味する。
	
	\subsection{電弱スケールと宇宙定数の間のフラクタル橋渡し}
	\label{subsec:bridge}
	
	もしシナリオ B が正しく，基本調整スケールが \(\LUV \approx 8.96\) TeV であるならば，宇宙論への深い繋がりが現れる。YUCT では，真空は単一の連続媒質ではなく，調整層のフラクタル階層である。スケール \(\LUV\) は標準模型の場が調整形状因子を経験する層を特徴付ける。しかし，より深い（赤外）層が存在し，それは次第に低いエネルギー尺度に対応する。宇宙定数，すなわち観測された暗黒エネルギー密度は，最も深い赤外層によって生成され，そのスケール \(\Lambda_{\mathrm{IR}}\) は，すべての量子場の零点エネルギーを完全なフラクタル階層にわたって積分したときに，観測値 \(\rho_{\Lambda} \approx 2.51 \times 10^{-47}\) GeV\(^4\) を再現するという要件によって設定される（完全な導出については付録 M を参照）。
	
	UV スケール \(\LUV\) と IR スケール \(\Lambda_{\mathrm{IR}}\) の関係は，YUCT の基本スケーリング量子によって支配される：
	\begin{equation}
		q = \left(\frac{3}{2}\right)^{1/3} \approx 1.1447.
	\end{equation}
	二つのスケールは整数 \(N\) 個の調整層によって隔てられている：
	\begin{equation}
		\LUV = \Lambda_{\mathrm{IR}} \cdot q^{N}.
		\label{eq:bridge}
	\end{equation}
	
	観測された宇宙定数から \(\Lambda_{\mathrm{IR}} \approx 0.0032\) eV が抽出される（付録 M，式 (M.47) 参照）。\(\LUV \approx 8.96 \times 10^{12}\) eV を用いると，
	\begin{equation}
		N = \frac{\ln(\LUV / \Lambda_{\mathrm{IR}})}{\ln q}
		= \frac{\ln(8.96 \times 10^{12} / 0.0032)}{\frac{1}{3}\ln(3/2)}
		\approx \frac{35.564}{0.135155} \approx 263.1.
		\label{eq:N_layers}
	\end{equation}
	を得る。
	
	最も近い半整数に丸めると（フェルミオン層に対する YUCT の代数的ループの要求），基本構造数
	\begin{equation}
		\boxed{N = 263.0}.
	\end{equation}
	が得られる。
	
	これは驚くべき結果である。電弱スケール（ヒッグスボソンの質量）と宇宙定数は，自由パラメータなしに，調整階層の正確に \(263\) 個の離散ステップによって結ばれている。宇宙定数の観測された小ささは謎ではない；それは単に素粒子物理学と宇宙論を隔てる調整ネットワークの巨大な深さを反映している。
	
	\textbf{解釈：} 階層の各層は「スケール変換器」として機能する：与えられた層の零点エネルギーは，次のより深い層の調整形状因子によって抑制される。全真空エネルギーはすべての層からの寄与の合計であり，各層は前の層よりも \(\beta = 2/3\) 倍小さい。幾何級数は有限値に収束し，それは正確に観測された暗黒エネルギー密度である（詳細な総和については付録 M を参照）。整数 \(N \approx 263\) は電弱スケールとニュートリノ質量スケール（最も深い物理的にアクセス可能な層）の間に収まる層の数である。この数が半整数として現れることは，YUCT フレームワークの非自明な整合性チェックであり，シナリオ B をさらに支持する。
	
	\subsection{二つのシナリオの比較}
	
	\begin{table}[H]
		\centering
		\caption{シナリオ A とシナリオ B の比較。}
		\label{tab:scenarios}
		\small
		\begin{tabular}{@{}lcc@{}}
			\toprule
			\textbf{性質} & \textbf{シナリオ A} & \textbf{シナリオ B} \\
			\midrule
			UV スケール \(\LUV\) & \(3\pi\Mpl \approx 10^{20}\) GeV & \(\approx 8.96\) TeV \\
			最小長 \(\Lmin\) & \(\frac{2}{3}\Lpl \approx 10^{-35}\) m & \(\approx 10^{-19}\) m \\
			ラグランジアン中の裸の質量 & 非ゼロ（微調整） & ゼロ（スケール不変性） \\
			放射補正 \(\delta m_H\) & \(10^{18}\) GeV & 125 GeV \\
			実験的アクセス可能性 & プランクスケール（アクセス不可） & LHC エネルギー（アクセス可） \\
			必要な微調整 & あり（\(10^{16}\) 分の１） & なし \\
			\bottomrule
		\end{tabular}
	\end{table}
	
	% ============================================================
	\section{シナリオ B の反証可能な予測}
	\label{sec:predictions}
	
	もしシナリオ B が正しいならば，以下の定量的予測を LHC および将来の衝突型加速器でテストできる：
	
	\begin{enumerate}[label=(\arabic*)]
		\item \textbf{ドレル–ヤン断面積の抑制}。測定されたドレル–ヤン断面積と標準模型断面積の比は，ダイレプトン不変質量 \(m_{\ell\ell}\) の関数として
		\[
		R(m_{\ell\ell}) \approx \exp\!\Bigl[-2\Bigl(\frac{m_{\ell\ell}^2}{\LUV^2}\Bigr)^{\!2/3}\,\Bigr],
		\]
		に従うべきである。ここで \(\LUV \approx 8.96\) TeV である。\(m_{\ell\ell} \approx 5.4\) TeV で 10\% の抑制が予測され，\(m_{\ell\ell} \approx 10\) TeV で \(\sim 2\) 倍に成長する。
		
		\item \textbf{ヒッグス対生成}。形状因子は高運動量移行でヒッグス自己結合を修正する。二ヒッグス不変質量分布は \(m_{HH} \gtrsim 3\) TeV で標準模型から逸脱するはずである。
		
		\item \textbf{フェルミオン質量}。もしすべてのフェルミオン質量も放射的に生成されるならば，それらの湯川結合は，放射補正が観測された質量スペクトルを再現するように選ばれなければならない。これはフェルミオン調整深度 \(L_f\) と共鳴因子 \(\xi_f\) を放射補正に関係付ける自己無撞着条件を課す。詳細な解析は将来の研究に委ねる。
	\end{enumerate}
	
	\subsection{微細構造定数に関する注記}
	
	トポロジカル不変量 \(N_{\mathrm{gauge}} = 12\)（付録 G）からの \(\alpha^{-1} \approx 137.036\) の導出は \(\LUV\) の値とは独立である。ゲージ結合は対数的に走る；\(\LUV\) での初期値は \(F_{18}\) のトポロジーによって固定される。シナリオ A では \(\LUV \sim \Mpl\) であり，プランクスケールから電弱スケールまでの対数走行は \(\alpha^{-1} \approx 137.036\) を与える。シナリオ B では \(\LUV \sim 9\) TeV であり，走行ははるかに短い区間にわたる。その場合，\(\alpha^{-1}\) のトポロジカル予測は TeV スケールで適用され，低エネルギー値は小さな対数発展によって得られる。これは \(\alpha^{-1}\) の数値予測を修正し，CODATA 値と照合されなければならない。予備的推定では，このずれはトポロジカル予測の現在の不確かさの範囲内にあることを示唆しているが，正確な計算が必要であり，付録 G の将来の改訂で提示される予定である。
	
	% ============================================================
	\section{議論：YUCT は万物の理論か，それとも電弱スケールの理論か？}
	\label{sec:discussion}
	
	元々の YUCT フレームワークは，プランク質量から電子質量までのすべての物理的スケールを単一の調整階層によって説明することを志していた。修正されたスカラー質量計算は，この野心の中の緊張を明らかにする：もし調整ネットワークがプランクスケールで動作するならば，放射補正は巨大であり，微調整された裸の質量によって打ち消されなければならない。
	
	シナリオ B は，YUCT 質量梯子を再解釈することによってこの緊張を解決する。深さ \(L\) を通じて質量を直接生成する代わりに，梯子は \textbf{湯川結合} を決定し，それが今度は \(\LUV \sim 9\) TeV のスケールで質量を放射的に生成する。プランクスケールはその後，調整ネットワークが \textit{形成され始める} スケールとなるが，物理的な粗さ – 形状因子が重要になるスケール – ははるかに低い。
	
	この再解釈は，微調整問題を排除しながら YUCT の代数的美しさを保持する。また，YUCT を投機的な量子重力理論から，現在および将来の衝突型加速器でテスト可能な反証可能な電弱物理のモデルへと変える。
	
	% ============================================================
	\section{結論}
	\label{sec:conclusion}
	
	我々は YUCT 内のスカラー質量補正における計算誤差を修正し，正確な解析値 \(C_0 = 3\sqrt{\pi}/(8\sqrt{2}) \approx 0.46999\) を得た。プランクスケール同一視 \(\LUV = 3\pi\Mpl\) を用いると，ヒッグス質量への放射補正は \(\delta m_H \approx 1.57\times 10^{18}\) GeV となり，裸の質量との微調整された打ち消しが必要となる。
	
	代替案として，我々はシナリオ B を提案した。そこではプランクスケールの裸のラグランジアンはスケール不変であり，すべての質量は放射的に生成される。自己無撞着性は，基本調整スケールが \(\LUV \approx 8.96\) TeV – LHC でアクセス可能な値 – であることを要求する。このシナリオは高エネルギー散乱過程に対する具体的で反証可能な予測を与える。
	
	シナリオ A（美的に不満足であるが YUCT のプランクスケール起源と整合する）とシナリオ B（予測力がありテスト可能であるが YUCT 質量梯子の再解釈を必要とする）の間の選択は，理論的な根拠だけでは行えない。それは実験によって決定されなければならない。我々は LHC 共同研究に，Eq.~(\ref{eq:F}) で \(\LUV \sim 9\) TeV として予測された形状因子抑制の証拠を探すために，高質量ドレル–ヤンおよびダイジェットデータを調べるよう促す。ヌル結果はシナリオ B を除外し，シナリオ A は論理的に可能なままであるが，それは階層性問題の解決としての YUCT の主張を弱めるであろう。
	
	\section*{謝辞}
	著者は YUCT セミナーの参加者との批判的議論に感謝する。
	
	\section*{データの利用可能性}
	すべての計算は自己完結している；数値評価に使用した Python コードは YPSDC リポジトリ \url{https://github.com/Alexey-Yakushev-YUCT/YPSDC} で入手可能である。
	
	\begin{thebibliography}{99}
		\bibitem{AppendixL} A.\,V.\,Yakushev, ``付録 L: フラクタル調整誤差スケーリング,'' in \emph{YUCT}, Zenodo (2026).
		\bibitem{AppendixV} A.\,V.\,Yakushev, ``付録 V: 調整プロセスのエントロピーとしての時間,'' in \emph{YUCT}, Zenodo (2026).
		\bibitem{AppendixG} A.\,V.\,Yakushev, ``付録 G: 調整ネットワークの幾何学的張力としての重力と YUCT における循環宇宙論,'' in \emph{YUCT}, Zenodo (2026).
		\bibitem{AppendixM} A.\,V.\,Yakushev, ``付録 M: YUCT 宇宙論 – 調整相転移としてのインフレーション,'' in \emph{YUCT}, Zenodo (2026).
		\bibitem{AppendixLambda} A.\,V.\,Yakushev, ``付録 \(\Lambda\): YUCT における階層性問題と宇宙定数問題の解決,'' in \emph{YUCT}, Zenodo (2026).
		\bibitem{AppendixJ} A.\,V.\,Yakushev, ``付録 J: YUCT トポロジカルオフセットからの標準模型フレーバーパラメータの完全な導出,'' in \emph{YUCT}, Zenodo (2026).
		\bibitem{AppendixFerra} A.\,V.\,Yakushev, ``付録 Ferra1.2: 秩序相と無秩序相のための普遍的調整定数,'' in \emph{YUCT}, Zenodo (2026).
	\end{thebibliography}
	
\end{document}